\section{Tái tạo lập thể đa góc nhìn:}
\label{sec:mvs}

Sau khi SfM đã cung cấp một tập hợp các hình ảnh $I = \{I_i\}$ cùng với camera pose $P = \{P_i\}$ và các tham số nội tại $K = \{K_i\}$ tương ứng, đồng thời tạo ra một đám mây điểm 3D thưa, mục tiêu của thuật toán Tái tạo lập thể đa góc nhìn (MVS) đó là tái tạo một cấu trúc hình học mây điểm 3D chi tiết, dày đặc hơn. Khác với SfM tập trung vào việc tìm các điểm đặc trưng nổi bật và ước tính các tham số máy ảnh, MVS hướng tới việc ước tính thông tin chiều sâu cho phần lớn các điểm ảnh cho các hình ảnh đầu vào, từ đó xây dựng nên một mô hình mây điểm 3D dày đặc của cảnh \cite{Seitz2006comparison, Furukawa2015multi}.

\subsection{Tương đồng quang học}
\label{subsec:mvs_principle}

Nguyên lý nền tảng, cốt lõi của hầu hết các thuật toán MVS là sự tương đồng quang học. Nguyên lý này dựa trên giả định rằng một điểm vật lý $X$ trong không gian 3D khi được quan sát từ các góc nhìn khác nhau (bởi các máy ảnh $P_i, P_j, ...$) sẽ có cùng sự tương đồng thị giác chẳng hạn như màu sắc hoặc cường độ sáng \cite{Seitz2006comparison}.

Nói cách khác, nếu chúng ta có một giả thuyết chính xác về vị trí 3D của một điểm ảnh $x = (u, v)$ trong ảnh tham chiếu $I_{ref}$, thì khi chiếu điểm 3D đó lên các ảnh quan sát ở góc nhìn khác, vùng ảnh xung quanh các điểm chiếu này phải có sự tương đồng cao về mặt quang học với vùng ảnh xung quanh $x$ trong $I_{ref}$. MVS khai thác nguyên lý này để tìm ra vị trí 3D thông qua độ sâu chính xác nhất cho mỗi điểm ảnh.

\subsection{Mô hình hóa bài toán và ước tính độ sâu}
\label{subsec:mvs_formulation}

Mục tiêu chính của MVS là ước tính bản đồ độ sâu $D_{ref}$ cho một ảnh tham chiếu $I_{ref}$. Mỗi giá trị $D_{ref}(x)$ tại điểm ảnh $x = (u,v)$ biểu diễn khoảng cách từ tâm máy ảnh $O_{ref}$ đến điểm 3D $X$ tương ứng với $p$. Điểm $X$ nằm trên tia nhìn $\mathbf{r}(d)$ xuất phát từ $O_{ref}$ và đi qua $x$. Tia nhìn này có thể được tham số hóa bởi độ sâu $d$:
\begin{equation}
    X(d) = \pi^{-1}(x, d, P_{ref}, K_{ref})
\end{equation}
trong đó $\pi^{-1}$ là hàm chiếu ngược từ điểm ảnh $x$ với độ sâu $d$ ra không gian 3D, sử dụng ma trận chiếu $P_{ref}$ và ma trận nội tại $K_{ref}$.

Để tìm độ sâu $d^*(x)$ tối ưu cho điểm ảnh $x$, MVS đánh giá một hàm chi phí hoặc hàm tương đồng dựa trên nguyên lý tương đồng quang học. Xét một tập các ảnh $\mathcal{N}(ref)$ lân cận với ảnh $I_{ref}$. Với mỗi độ sâu giả định $d_h$ dọc theo tia nhìn của $x$, điểm 3D tương ứng $X(d_h)$ được chiếu lên từng ảnh nguồn $I_k$ ($k \in \mathcal{N}(ref)$) để có điểm ảnh $x'_k = \pi(X(d_h), P_k, K_k)$, với $\pi$ là hàm chiếu 3D-2D. Sau đó, sự tương đồng quang học giữa một vùng ảnh (gọi là patch) $W$ xung quanh $x$ trong $I_{ref}$ và các vùng ảnh xung quanh $x'_k$ trong $I_k$ được đo lường.

Một hàm chi phí tương đồng quang học $C(x, d_h)$ phổ biến là tổng bình phương độ lệch hoặc tổng trị tuyệt đối độ lệch giữa các vùng ảnh:
\begin{equation}
    C_{SSD}(x, d_h) = \sum_{k \in \mathcal{N}(ref)} \sum_{\Delta x \in W} \| I_{ref}(x+\Delta x) - I_k(x'_k+\Delta x) \|^2
\end{equation}

Một độ đo tương đồng phổ biến khác, có ưu điểm là mạnh mẽ hơn đối với sự thay đổi tuyến tính về độ sáng và độ tương phản giữa các ảnh là NCC (Normalized Cross-Correlation - Tương quan chéo đã chuẩn hóa). Với vùng ảnh $W_{ref}$ xung quanh điểm ảnh $x$ trong ảnh tham chiếu $I_{ref}$ và vùng ảnh $W_k$ xung quanh điểm ảnh $x'_k$ (hình chiếu của điểm 3D ứng với độ sâu $d_h$) trong ảnh nguồn $I_k$, giá trị NCC được tính như sau:
\begin{equation} \label{eq:ncc}
    \text{NCC}(W_{ref}, W_k) = \frac{\sum_{u,v \in W} (I_{ref}(u,v) - \bar{I}_{ref}) (I_k(u,v) - \bar{I}_{k})}{\sqrt{\sum_{u,v \in W} (I_{ref}(u,v) - \bar{I}_{ref})^2 \sum_{u,v \in W} (I_k(u,v) - \bar{I}_{k})^2}}
\end{equation}
trong đó $\bar{I}_{ref}$ và $\bar{I}_{k}$ lần lượt là giá trị cường độ trung bình của các điểm ảnh trong vùng ảnh $W_{ref}$ và $W_k$, $I_{ref}(u,v)$ và $I_k(u,v)$ là giá trị cường độ tại tọa độ $(u,v)$ tương đối bên trong các vùng ảnh, và tổng $\sum_{u,v \in W}$ được lấy trên tất cả các điểm ảnh $(u,v)$ trong vùng ảnh $W$.

Giải thích công thức NCC:
\begin{itemize}
    \item Tử số: 
    \begin{equation}
        \sum_{u,v \in W} (I_{ref}(u,v) - \bar{I}_{ref}) (I_k(u,v) - \bar{I}_{k})    
    \end{equation}
    Đây là hiệp phương sai giữa hai vùng ảnh \(W_{ref}\) và \(W_k\), đo lường mức độ biến động chung của cường độ pixel sau khi đã trừ giá trị trung bình (\(\bar{I}_{ref}\) và \(\bar{I}_{k}\)). Việc trừ đi giá trị trung bình nhằm loại bỏ ảnh hưởng của cường độ ánh sáng tới việc tính toán sự tương đồng mà chỉ tập trung vào sự thay đổi tương đối của cường độ điểm ảnh.

    \item Mẫu số:
    \begin{equation}
        \sqrt{\sum_{u,v \in W} (I_{ref}(u,v) - \bar{I}_{ref})^2 \sum_{u,v \in W_k} (I_k(u,v) - \bar{I}_{k})^2}
    \end{equation}
    Đây là tích của độ lệch chuẩn (hoặc căn bậc hai của phương sai) của hai vùng ảnh. Cụ thể:
    \begin{itemize}
        \item \(\sum_{u,v \in W} (I_{ref}(u,v) - \bar{I}_{ref})^2\) là tổng bình phương độ lệch của cường độ điểm ảnh trong \(W_{ref}\), tương ứng với phương sai (chưa chia cho số điểm ảnh).
        \item Tương tự cho \(W_k\).
        \item Căn bậc hai của tích hai tổng này chuẩn hóa hiệp phương sai, đảm bảo giá trị NCC nằm trong khoảng \([-1, 1]\).
    \end{itemize}

    \item Tính chất chuẩn hóa: Công thức này đảm bảo NCC không bị ảnh hưởng bởi sự thay đổi tuyến tính về độ sáng hoặc độ tương phản, vì các giá trị cường độ đã được trừ đi trung bình và chuẩn hóa bởi độ lệch chuẩn.
\end{itemize}

Giá trị NCC nằm trong khoảng $[-1, 1]$, với giá trị gần $1$ biểu thị sự tương đồng cao nhất (hoàn toàn tương quan tuyến tính dương), giá trị gần $-1$ biểu thị sự đối nghịch cao nhất (tương quan tuyến tính âm), và giá trị gần $0$ biểu thị không có tương quan tuyến tính. Do đó, khi sử dụng NCC làm độ đo sự tương đồng quang học, mục tiêu là tìm độ sâu $d^*(x)$ để tối đa hóa giá trị tương đồng NCC thay vì tối thiểu hóa chi phí như SSD hay SAD:
\begin{equation} \label{eq:ncc_maximization}
    d^*(x) = \underset{d_h}{\arg\max} \, \left( \sum_{k \in \mathcal{N}(ref)} w_k \cdot \text{NCC}(W_{ref}, W_k(d_h)) \right)
\end{equation}
Trong đó $W_k(d_h)$ là patch trong ảnh nguồn $I_k$ tương ứng với độ sâu giả định $d_h$, và $w_k$ là trọng số tùy chọn cho từng ảnh nguồn $k$.

\subsection{Các bước thực hiện thuật toán}
\label{subsec:mvs_steps}

Một quy trình MVS điển hình, ví dụ như các phương pháp dựa trên vùng ảnh (gọi là patch-based), thường bao gồm các bước sau:

\begin{itemize}
    \item Chọn cặp ảnh tham chiếu - nguồn: Với mỗi ảnh có thể dùng làm ảnh tham chiếu $I_{ref}$, lựa chọn một tập các ảnh nguồn $I_{src}$ phù hợp, có đủ sự chồng lấp nhưng cũng có đủ đường cơ sở để đảm bảo độ chính xác cho phép giao hội, phép chiếu.
    \item Ước tính bản đồ độ sâu: Với mỗi ảnh tham chiếu $I_{ref}$, thực hiện tìm kiếm độ sâu $d^*(x)$ tối ưu cho từng điểm ảnh $x$ bằng cách tối ưu hóa hàm chi phí tương đồng quang học với các vùng chiếu trên các ảnh nguồn tương ứng. Kết quả là một bản đồ độ sâu $D_{ref}$ và có thể cả bản đồ pháp tuyến $N_{ref}$ cho mỗi ảnh tham chiếu.

    \item Lọc bản đồ độ sâu: Các bản đồ độ sâu thô thường chứa nhiễu và các giá trị ngoại lai. Các bộ lọc được áp dụng để loại bỏ các ước tính không đáng tin cậy, ví dụ bằng cách kiểm tra sự nhất quán hình học giữa các bản đồ độ sâu của các view nhìn thấy cùng một vùng ví dụ: điểm 3D ước tính từ $I_{ref}$ phải chiếu gần với điểm 3D ước tính từ $I_{src}$ khi $I_{src}$ làm tham chiếu.
    \item Hợp nhất bản đồ độ sâu: Các điểm 3D từ các bản đồ độ sâu đã được lọc được chiếu vào một không gian 3D chung. Các điểm từ các góc nhìn khác nhau mô tả cùng một bề mặt được hợp nhất lại, thường thông qua các cấu trúc như lưới voxel hoặc trực tiếp tính trung bình có trọng số, để tạo thành một đám mây điểm 3D dày đặc cuối cùng. Các điểm ngoại lai còn sót lại của mô hình mây điểm có thể được loại bỏ bằng các thuật toán hậu xử lý \cite{schoenberger2016mvs}.
\end{itemize}

Ưu điểm của thuật toán MVS truyền thống:
\begin{itemize}
    \item Có khả năng tạo ra mô hình 3D rất chi tiết và có độ chính xác hình học cao, đặc biệt với các cảnh có đặc trưng phong phú và ổn định.
    \item Không yêu cầu dữ liệu huấn luyện lớn như các phương pháp học sâu.
    \item Các thuật toán và lý thuyết đã được phát triển và kiểm chứng trong thời gian dài.
\end{itemize}

Nhược điểm của MVS truyền thống:
\begin{itemize}
    \item Rất nhạy cảm với độ chính xác của camera pose đầu vào từ SfM. Khi camera pose sai lệch sẽ dẫn đến kết quả MVS sai lệch nghiêm trọng, đám mây điểm có độ chính xác cấu trúc hình học kém.
    \item Gặp nhiều khó khăn với các cảnh, đối tượng không có hoặc ít đặc trưng bề mặt, bề mặt phản chiếu hoặc trong suốt, do nguyên lý sự tương đồng quang học khó áp dụng và không phù hợp với các trường hợp này.
    \item Nhạy cảm với sự che khuất và thay đổi điều kiện chiếu sáng mạnh giữa các ảnh.
    \item Yêu cầu thời gian và tài nguyên tính toán đáng kể khi tập ảnh có ảnh độ phân giải cao và số lượng ảnh lớn.
\end{itemize}
Những hạn chế này, đặc biệt là sự khó khăn khi giải quyết các đối tượng, cảnh với bề mặt thiếu đặc trưng, là động lực để xem xét các phương pháp cải tiến hoặc các hướng tiếp cận thay thế như NeRF cho các đối tượng thách thức như cột ăng-ten.